\scanpage{60}
\mspage{60}
\hfill n°269

$K$-\uline{espace vectoriel des formes sesquilinéaires sur} $E\times E'$. \uline{Si l'on a de plus} $E=E'$, \uline{pour que} $\Phi$ \uline{soit hermitienne, il faut et il suffit que} $\Phi_1$ \uline{soit symétrique} (resp. \uline{que} $\Phi_2$ \uline{soit antisymétrique}).
\margin{\unsure{''}}

En effet toute application $\Phi$ de $E\times E'$ dans $A$ s'écrit, d'une manière et d'une seule, sous la forme $\Phi=\Phi_1+i\Phi_2$, où $\Phi_1$ et $\Phi_2$ sont des applications de $E\times E'$ dans $K$. Pour que l'application partielle $x\longrightarrow\Phi(x,x')$ soit $A$-linéaire, il faut et il suffit, d'après la formule (\unsure{5|3}) (resp. (4)) du n°2, que $\Phi_1$ (resp. $\Phi_2$) soit $K$-linéaire en $x$ et que l'on ait
\[
(9)\qquad\Phi(x,x')=\Phi_1(x,x')-i\Phi_1(j(x),x')
\]
\[
(\text{resp.}\quad(10)\qquad\Phi(x,x')=\Phi_2(j(x),x')+i\Phi_2(x,x')).
\]
De même, pour que $\overline{\Phi(x,x')}$ soit $A$-linéaire en $x'$, il faut et il suffit que $\Phi_1$ (resp. $\Phi_2$) soit $K$-linéaire en $x'$ et que l'on ait
\[
(11)\qquad\Phi(x,x')=\Phi_1(x,x')+i\Phi_1(x,j'(x'))
\]
\[
(\text{resp.}\quad(12)\qquad\Phi(x,x')=\Phi_2(x,j'(x'))+i\Phi_2(x,x')\ ).
\]
Il en résulte immédiatement que, pour que $\Phi$ soit sesquilinéaire, il faut et il suffit qu'elle s'écrive sous l'une ou l'autre des formes (9) et (11) (resp. 10) et (12)) avec $\Phi_1$ (resp. $\Phi_2$) $K$-bilinéaire invariante par $j$ et $j'$.

Il résulte de ceci que, pour qu'une forme sesquilinéaire $\Phi=\Phi_1+i\Phi_2$ soit nulle, il faut et il suffit que $\Phi_1$ (resp. $\Phi_2$) soit nulle. Or, si $E=E'$, on a
\[
\Phi(y,x)=\Phi_1(y,x)+i\Phi_2(y,x)\quad\text{et}\quad\overline{\Phi(x,y)}=\Phi_1(x,y)-i\Phi_2(x,y)\ ;
\]
pour que ces deux expressions soient égales, autrement dit pour
\margin{59}
